% ch2 C.5 有效哈密顿量理论的严格基础 (原书页 91–95 / PDF p106–p110)

% 衔接说明：原书 p.91（PDF p106）小节标题之上尚有 C.4（“击穿”效应）的收尾一段
% （“It is worth noting that in this limit ... asymptotically convergent at best.”），
% 按分工自本小节标题起译，该段完全略去。页首的 Figure 25 图身落在 p106，其正文引用
% （“……便可以跳入下一条轨道（图 25）”）已见于 ch2_C34.tex，%FIG 标记置于此处。
% 本小节正文另引 Figure 26、27、28，均出现在本文件范围内，随文标记。
% 本文件末句在原书 p.95（PDF p110）以句号完整结束；p.96（PDF p111）起为第 3 章。

% ---- 原书 p.91 (PDF p106) ----

%FIG{25}: 图 25　极开放晶格中紧束缚 Bloch 波函数的实部示意：波函数逐格点平滑延伸，绕过格点处经虚线所示的一次跳跃（Hops）过渡（原书图注仅作 “Figure 25”）

\subsection{有效哈密顿量理论的严格基础}

现在我想简要讨论一下，人们如何着手改进有效哈密顿量理论 (56)。Adams 提出的一个非常简单的定性图像，将以直观的图形方式展示其中的原理。考虑一个非常开放的晶格，其中紧束缚理论是相当好的近似，于是 Bloch 函数就是紧束缚函数，
\begin{equation*}
b^n_k = \sum_j e^{ik\cdot R_j}\,\phi_n(r - R_j)
\end{equation*}
% ---- 原书 p.92 (PDF p107) ----
我们假定交叠相当小，因而在相当好的近似下，$\phi_n$ 也就是 Wannier 函数 $a_n(r - R_j)$。Bloch 波函数实部的图像便可能如图 26 所示。

%FIG{26}: 图 26　紧束缚 Bloch 波函数实部示意：缓变的 $b_n(k)$ 乘以各格点上的 Wannier 函数 $a_n(r - R_j)$（钟形峰），过零后峰反号；横轴为 $r$（原书图注仅作 “Figure 26”）

现在，在推导有效哈密顿量时我们做了一个基本假设：$V(r)$——内部场产生的真实微扰——可以用 $V(R)$ 代替，其中 $V(R)$ 由如下关系定义：
\begin{equation*}
V(R)\,a_n(r - R_j) = V(R_j)\,a_n(r - R_j)
\end{equation*}
例如，对均匀电场而言，这意味着 $r$ 可以用相应 Wannier 函数中心处的值代替。于是很容易看出，在我们正在讨论的高度局域情形中，如果 $V(r) = -er$ 是图中所示的线性势，那么 $V(R)$ 就是图 27 所示的阶梯势，而二者之差就是图 28 所示的锯齿波。让我们把这个差记为 $X = r - R$。显然，在任何实际情形中 $X$ 都不会长得像图中那样，但它总可以用这个差来定义，而这个差又依赖于 $a_n(r - R_j)$。我们立刻注意到：

1. 差 $X$ 是 $eEa$ 量级的小量，其中 $a$ 是晶格常数，于是微扰的“大”完全包含在算符 $R$ 之中，从而有效哈密顿量理论本质上包含了所有\emph{大的}效应。

2. $X$ 以晶格周期为周期。因此，在微扰论的任意阶上，我们都可以把 Bloch 函数和 Wannier 函数重新定义为如下\emph{新}能带问题的解：
\begin{equation*}
\mathcal{H} b_k = \left(\frac{p^2}{2m} + V(r) - eEX\right) b_k \simeq {(E^n_k)}'\, b_k
\end{equation*}
% ---- 原书 p.93 (PDF p108) ----
%FIG{27}: 图 27　强局域情形下的线性势与阶梯势：自上而下依次为各格点上的 Wannier 函数 $a_n$（钟形峰）、线性势 $r$（斜直线）与阶梯下降的 $R$（平台之间以虚线相连），竖直线标记各格点位置（原书图注仅作 “Figure 27”）

%FIG{28}: 图 28　差 $X = r - R$ 随 $r$ 变化的锯齿波（原书图注仅作 “Figure 28”）

从而消去 $X$ 的效应。这样，在微扰论的任意阶上都存在一组新能带，带有新能量 ${(E^n_k)}'$、新速度（不再严格地由 $\partial E/\partial k$ 给出）等等，而有效哈密顿量理论在附带这一条件的情况下对新能带仍然有效。但要注意，$X$ 是用 $a$ \emph{定义}出来的，而 $a$ 又依赖于 $b_k$，因此求解该方程的唯一办法不是一劳永逸，而是对所有阶逐次迭代。

3. 这种情形下 $X$ 对 Wannier 函数的影响实际上是一个显而易见的物理效应，它确实被我们原有的理论遗漏了：电场不仅必定有加速电子的效应，它还有使单个原子的波函数极化的效应。例如，它必定会把一些 p 带函数混入 s 带中，等等；而且无论 Wannier 函数局域到什么程度，这一效应都必定存在。换句话说，十分明显，$X$ 可以具有带间矩阵元 $X^{nn'}$；但我们预期这些带间矩阵元将是某个周期势的矩阵元，是 $eEa$ 量级的小量。

% ---- 原书 p.94 (PDF p109) ----

讲到这里，这仅仅是一个纲领，而这些结果的实际证明，正是我前面提到的那一组人在过去 10 年里所从事的工作。为了看清必须做些什么，让我们离开紧束缚近似，对一个一般能带来定义 $X$。实际上有两种办法；最显然的一种是用 Wannier 函数来表述：
\begin{equation*}
X\,a_n(r - R_j) = (r - R_j)\,a_n(r - R_j)
\end{equation*}
于是 $X$ 的矩阵元就是
\begin{equation*}
(n, j\,|\,X\,|\,n', j') = \int dr\; a_n(r - R_j)\,a_{n'}(r - R_{j'})\,(r - R_{j'})
\end{equation*}
因此 $X$ 是 $R_j - R_{j'}$ 的函数，这就证明了它是周期的，但非局域。

不那么显然的办法是所谓“晶体动量表象”（crystal momentum representation）。显然，由于
\begin{equation*}
b^n_k(r) = \sum_j e^{ik\cdot R_j}\,a_n(r - R_j)
\end{equation*}
便有
\begin{equation*}
R\,b^n_k = -\,i\,\frac{\partial}{\partial k}\,b^n_k
\end{equation*}
而用另一种表象，
% 校注：按链式法则，此式末项应为 $-i\,e^{ik\cdot r}\frac{\partial}{\partial k}u_k(r)$；
% 原书末项将因子 $-i$ 印漏，此处照原书转录。
\begin{equation*}
-\,i\,\frac{\partial}{\partial k}\,b^n_k(r) = r\,e^{ik\cdot r}\,u_k(r) - e^{ik\cdot r}\,\frac{\partial}{\partial k}\,u_k(r)
\end{equation*}
所以 $X$ 就是单独作用在周期部分 $u_k$ 上的运算 $-i(\partial/\partial k)$。同样，这显然是周期的，因为结果是一个具有同一 $k$ 的 Bloch 函数；但算符是非局域的。这种晶体动量表象在磁场情形中更有价值，因为在那里同样有必要把微扰中出现的算符 $p$ 展开成周期与非周期两部分。

在这两种表象中，每一种都出现一个可能的困难。第一种表象中最为明显：如果 $a_n(r - R_j)$ 局域得不够充分，$X$ 的矩阵元可能很大甚至发散。例如，Wannier 在其原始论文中计算的情形是把自由电子当作能带来考虑的情形，$a_n$ 仅按 $\sin r/r$ 衰减，因此矩阵元会发散得相当厉害。

绕开这一困难的办法是 Gibson 与 Blount 的贡献。注意我们定义了
% ---- 原书 p.95 (PDF p110) ----
\begin{equation*}
a_n(r - R_j) = \sum_k e^{ik\cdot R_j}\,b^n_k(r)
\end{equation*}
然而，这绝不是唯一的，因为各个 $b_k$ 彼此之间可以带有完全任意的相位，它们只是在相差一个任意常数因子的意义上被定义的。于是我们拥有一个由可能的 $a_n$ 组成的无穷空间：
\begin{equation*}
a_n = \sum_k e^{ik\cdot R_j}\,e^{i\phi_k}\,b^n_k(r)
\end{equation*}
正如 Blount 所证明的，这种相位变换非常像一个规范变换（gauge transformation），多数结果并不因它而真正改变；但方便的做法是以 $a_n$ \emph{尽可能局域}为要求来规定各相位——也就是说，通过一个变分原理。

在 $\partial/\partial k$ 的定义中我们可以看到完全同样的麻烦；同样，每个 $b_n$ 都可以乘以一个任意相位，只有凭借某种特定的选取，我们才能使 $X$ 成为 $k$ 的光滑函数。

在单一、分离的能带情形，可以证明在最佳的相位选取下 $a_n$ 随距离指数式衰减，整个理论收敛得非常好。当出现简并线或简并点时，$a_n$ 衰减得不太快，不过事实上它们衰减得恰好足够快，足以保证收敛。但在这样的点或线上，通过在不同能带之间作新的线性组合，使 Wannier 函数良好局域、$X$ 与 $k$ 算符行为良好，便可以得到简单得多的结果。然而，这时未微扰的能带能量不再是一个数值函数 $E_n(k)$，而是一个矩阵 $E_{nn'}(k)$。这一做法首先由 Luttinger 在处理 Ge 中简并价带顶处的回旋共振（cyclotron resonance）时使用。利用这种相位选取技术，Blount 已经证明了有效哈密顿量理论对微扰论所有阶的有效性。

最后，不妨不再深入细节，对由这些结果可以得到的那种相当迷人的复杂性略作评论。一个重要结果是如下事实：在 $E$ 或 $H$ 的各个阶上，各种物理算符中会进入新的、相当出人意料的项——例如，反常 Hall 效应（anomalous Hall effect）(57) 中著名的反常电流——它可以导致电流、能量输运等等，这些纯粹是量子极化效应，处于正常输运理论的领域之外。当诸如自旋-轨道耦合（spin-orbit coupling）这类进一步的复杂性被恰当地引入时，还会出现许多其他效应，它们复杂但对专家来说饶有趣味。
